% Kompilacja: pdflatex prezentacja.tex (dwukrotnie).
% Motyw "Metropolis" odtworzony samodzielnie w preambule (pakiet beamertheme-metropolis nie jest
% zainstalowany w tym środowisku). Aby użyć oryginalnego motywu: tlmgr install beamertheme-metropolis
% pgfopts, a następnie zastąpić blok "MOTYW" linią \usetheme{metropolis}.
\documentclass[11pt,aspectratio=169]{beamer}
\usepackage[utf8]{inputenc}
\usepackage[T1]{fontenc}
\usepackage[polish]{babel}
\usepackage{lmodern}
\usepackage{amsmath,amssymb}
\usepackage{booktabs}
\usepackage{tikz}

% ======================= MOTYW w stylu Metropolis =======================
\usetheme{default}
\setbeamertemplate{navigation symbols}{}
\definecolor{mteal}{HTML}{2A9D8F}
\definecolor{morange}{HTML}{E07A2F}
\definecolor{mdark}{HTML}{2D3142}
\setbeamercolor{normal text}{fg=mdark}
\setbeamercolor{structure}{fg=mteal}
\setbeamercolor{title}{fg=morange}
\setbeamercolor{frametitle}{fg=morange,bg=white}
\setbeamercolor{itemize item}{fg=mteal}
\setbeamercolor{itemize subitem}{fg=mteal}
\setbeamerfont{title}{series=\bfseries,size=\LARGE}
\setbeamerfont{frametitle}{series=\bfseries,size=\Large}
\setbeamertemplate{itemize item}{\small\raise0.5pt\hbox{\textbullet}}
\setbeamertemplate{itemize subitem}{--}
% frametitle: tytuł u góry + cienka turkusowa linia
\setbeamertemplate{frametitle}{%
  \vspace*{0.7em}\hspace*{0.1em}%
  {\usebeamerfont{frametitle}\usebeamercolor[fg]{frametitle}\insertframetitle}\par
  \vspace*{0.15em}{\color{mteal}\rule{\textwidth}{1.1pt}}\par}
% footline: numer slajdu po prawej, cienka linia
\setbeamertemplate{footline}{%
  \hbox{\color{mteal}\rule{\paperwidth}{1pt}}%
  \hfill{\footnotesize\color{mdark!70}\insertframenumber\,/\,\inserttotalframenumber\ \ }\vspace*{2pt}}
% ========================================================================

% pomocniczy rysunek: dwa łuki nad punktami 1..4
\newcommand{\luki}[4]{%
\begin{tikzpicture}[baseline=-2pt,scale=0.6]
  \foreach \x in {1,2,3,4} { \fill (\x,0) circle (1.7pt); \node[below=1pt,font=\scriptsize] at (\x,0) {\x}; }
  \draw[red!80!black,very thick] (#1,0) to[bend left=70] (#2,0);
  \draw[blue!75!black,very thick] (#3,0) to[bend left=70] (#4,0);
\end{tikzpicture}}

\title{Arc Match}
\subtitle{Gra Maker--Breaker na uporządkowanych skojarzeniach}
\author{Wojciech Jasiński}
\date{Czerwiec 2026}

\begin{document}

\begin{frame}[plain]
  \titlepage
\end{frame}

\begin{frame}{Uporządkowane skojarzenia i wzorce}
  \begin{block}{Definicja}
    \emph{Uporządkowane skojarzenie} rozmiaru $n$ to $n$ łuków bez wspólnych końców nad $2n$ punktami.
    Dwa łuki są w jednym z trzech położeń:
  \end{block}
  \begin{center}
    \begin{tabular}{ccc}
      \luki{1}{2}{3}{4} & \luki{1}{4}{2}{3} & \luki{1}{3}{2}{4}\\[2pt]
      \textbf{ułożenie} & \textbf{zagnieżdżenie} & \textbf{krzyżowanie}
    \end{tabular}
  \end{center}
  Podskojarzenie jest \alert{jednorodne}, gdy \emph{każda} para łuków ma ten sam typ. \textbf{Przykład}
  ($n=3$): $\{(1,4),(2,5),(3,6)\}$ to 3-krzyżowanie, $\{(1,6),(2,5),(3,4)\}$ to 3-zagnieżdżenie,
  $\{(1,2),(3,4),(5,6)\}$ to 3-ułożenie.
\end{frame}

\begin{frame}{Zasady gry}
  \begin{itemize}
    \item Plansza: $2n$ punktów na prostej; parametry $n$, $k$, typ celu (domyślnie \emph{krzyżowanie}).
    \item Gracz w swoim ruchu łączy dwa nieużyte punkty łukiem w swoim kolorze.
    \item \alert{Maker (czerwony)} chce zbudować $k$ własnych łuków parami się krzyżujących.
    \item \alert{Breaker (niebieski)} stara się mu przeszkodzić.
    \item Maker rozpoczyna. Maker wygrywa w chwili powstania wzorca; Breaker --- gdy plansza się
      zapełni bez wzorca. Remis: na korzyść Breakera.
  \end{itemize}
\end{frame}

\begin{frame}{Twierdzenie i heurystyka}
  \begin{block}{Erdős--Szekeres dla skojarzeń (Dudek, Grytczuk, Ruciński, 2024)}
    Każde uporządkowane skojarzenie rozmiaru $n$ zawiera wzorzec jednorodny rozmiaru $\ge n^{1/3}$.
    Dotyczy to \emph{całego} skojarzenia (obu kolorów), więc \emph{nie} daje wprost dolnego
    ograniczenia dla Makera.
  \end{block}
  \begin{block}{Erdős--Selfridge --- tylko heurystyka}
    Gdyby łuki były niezależne, warunek $\binom{2n}{2k}2^{-k}<\tfrac12$ gwarantowałby wygraną
    Breakera. Lecz wybór łuku zużywa jego końce (Breaker może unieważnić cel, \emph{nie} zajmując go),
    więc to nie jest model Arc Match --- ES służy jako heurystyka, nie twierdzenie.
  \end{block}
\end{frame}

\begin{frame}{Zaimplementowane strategie}
  \begin{itemize}
    \item \textbf{Poziom 0 --- losowy}: dowolny legalny łuk.
    \item \textbf{Poziom 1 --- zachłanny}: maksymalizuje najdłuższy łańcuch danego typu, blokuje
      najgroźniejszy łańcuch przeciwnika.
    \item \textbf{Poziom 2 --- Erdős--Selfridge}: gra według potencjału
      $\Phi=\sum_{A\ \text{żywe}} 2^{-(k-r_A)}$.
    \item \textbf{Poziom 3 --- solver}: dokładny minimax (alpha--beta, tablica transpozycji,
      symetria) dla małych $n$.
  \end{itemize}
\end{frame}

\begin{frame}{Wyniki: skuteczność a poziom AI}
  Odsetek zwycięstw Makera (cel: krzyżowanie, $k=2$, Breaker na poziomie~2, 300 gier):
  \begin{center}
    \begin{tabular}{l|cccccc}
      \toprule
      $n$ & 3 & 4 & 5 & 6 & 7 & 8\\
      \midrule
      Maker poz.~0 (losowy) & 0\% & 6\% & 13\% & 26\% & 39\% & 53\%\\
      Maker poz.~2 (ES) & 0\% & 100\% & 100\% & 100\% & 100\% & 100\%\\
      \bottomrule
    \end{tabular}
  \end{center}
  \medskip
  Mocny Maker (poziom~2) wygrywa wszystkie 300 gier już od $n=4$, losowy --- dopiero stopniowo:
  poziom AI ma duże znaczenie. Że to rzeczywiście próg gry, potwierdza dokładny solver (następny
  slajd).
\end{frame}

\begin{frame}{Wyniki: próg gry $k^*(n)$}
  Dokładny próg minimax $k^*(n)$ --- największe $k\ge 2$, które Maker wymusza przy grze optymalnej
  (solver poziomu~3), cel: krzyżowanie:
  \begin{center}
    \begin{tabular}{l|cccccc}
      \toprule
      $n$ & 3 & 4 & 5 & 6 & 7 & 8\\
      \midrule
      $k^*(n)$ & 0 & 2 & 2 & 2 & 3 & 3\\
      pułap $\lceil n/2\rceil$ & 2 & 2 & 3 & 3 & 4 & 4\\
      $n^{1/3}$ & 1.44 & 1.59 & 1.71 & 1.82 & 1.91 & 2.00\\
      \bottomrule
    \end{tabular}
  \end{center}
  \medskip
  Dla $n\ge 4$ próg $k^*(n)$ leży poniżej pułapu $\lceil n/2\rceil$ i powyżej krzywej odniesienia
  $n^{1/3}$ (przy $n=3$ wzorzec $k=2$ jest możliwy, lecz Maker nie potrafi go wymusić). Ogólne
  zachowanie $k^*(n)$ dla dużych $n$ pozostaje \alert{problemem otwartym}.
\end{frame}

\begin{frame}[plain]
  \centering\vfill
  {\usebeamerfont{frametitle}\color{morange}Dziękujemy za uwagę}\\[0.6em]
  {\color{mteal}\rule{0.4\textwidth}{1.3pt}}
  \vfill
\end{frame}

\end{document}
